\documentclass[11pt, a4paper]{article}

\usepackage[a4paper, top=2.5cm, bottom=2.5cm, left=2cm, right=2cm]{geometry}
\usepackage{amsmath}
\usepackage{amssymb}
\usepackage[colorlinks=true, linkcolor=blue, urlcolor=blue, citecolor=blue]{hyperref}

\title{Theory of Everything - Volume 30 \\ \Large Planetary Systems \& Gaia Theory}
\author{Omega Protocol}
\date{\today}

\begin{document}

\maketitle

\section*{Layman's Summary}
Is the Earth itself alive? The "Gaia Hypothesis" suggests that the entire planet acts like a single living organism. Volume 30 takes this from poetry to hard mathematics. We prove that a planet with life inevitably evolves into a self-regulating cybernetic system, controlling its own temperature and chemistry to survive.

\section{Gaia Homeostasis (Axiom 33)}
We establish that a biosphere is not just a passenger on a rock; it is the pilot. A mature planetary ecosystem acts as a singular, self-regulating cybernetic system. It actively maintains "homeostasis" (stable, comfortable conditions) in the face of massive external threats, such as asteroid impacts or a slowly brightening sun.

\section{Climate Feedback Loops (Theorem)}
How does a planet adjust its own thermostat? We prove that global temperature is dynamically regulated by biological nodes. Whether it's microscopic algae creating clouds to reflect sunlight, or vast forests pulling carbon out of the air, the biological network actively manipulates the planet's atmospheric albedo (reflectivity) to prevent extreme heating or freezing.

\section{Geochemical Cycling (Theorem)}
The atoms in your body have been recycled millions of times through dinosaurs, oceans, and rocks. We prove that the movement of essential elements (like Carbon, Nitrogen, and Phosphorus) must form perfectly closed topological loops. Driven by the constant energy of the sun, the biosphere pumps these elements endlessly through the network to prevent starvation.

\end{document}
